
\documentclass{article}
\usepackage{colortbl}
\usepackage{makecell}
\usepackage{multirow}
\usepackage{supertabular}

\begin{document}

\newcounter{utterance}

\centering \large Interaction Transcript for game `cladder', experiment `full\_v1.5\_default', episode 7000 with qwen3.5{-}4b.
\vspace{24pt}

{ \footnotesize  \setcounter{utterance}{1}
\setlength{\tabcolsep}{0pt}
\begin{supertabular}{c@{$\;$}|p{.15\linewidth}@{}p{.15\linewidth}p{.15\linewidth}p{.15\linewidth}p{.15\linewidth}p{.15\linewidth}}
    \# & \multicolumn{2}{c}{Player} && \multicolumn{2}{c}{Game Master} \\
    \hline

    \theutterance \stepcounter{utterance}  
    & & & \multicolumn{4}{p{0.6\linewidth}}{
        \cellcolor[rgb]{0.9,0.9,0.9}{
            \makecell[{{p{\linewidth}}}]{
                \texttt{\tiny{[P1$\langle$GM]}}
                \texttt{You are an expert in causal inference. The following question is not a typical commonsense query, but rather a meticulously designed question created by a professor specializing in causal inference, intended to assess the students' mastery of the course content.} \\
\texttt{Imagine a self{-}contained, hypothetical world with only the following conditions, and without any unmentioned factors or causal relationships: Jyka has a direct effect on gyzp and hwax. Hwax has a direct effect on lirg. Gyzp has a direct effect on lirg. For those who are not jyka, the probability of lirg is 63\%. For those who are jyka, the probability of lirg is 49\%. For those who are jyka, would it be more likely to see lirg if the individual was not jyka?} \\
\texttt{Start your answer with 'Yes' or 'No', followed by additional reasoning or evidence to support your explanation.} \\
            }
        }
    }
     \\ \\

    \theutterance \stepcounter{utterance}  
    & \multicolumn{4}{p{0.6\linewidth}}{
        \cellcolor[rgb]{0.561,0.737,0.561}{
            \makecell[{{p{\linewidth}}}]{
                \texttt{\tiny{[P1$\rangle$GM]}}
                \texttt{No} \\
\\ 
\texttt{**Reasoning:**} \\
\\ 
\texttt{The question asks whether, for an individual who is **Jyka**, it would be more likely to see **Lirg** if that individual was **not** Jyka. This is a counterfactual query asking to compare the probability of \$Lirg\$ under two different conditions for the same subject:} \\
\texttt{1.  The actual condition: The individual is Jyka (\$Jyka = 1\$).} \\
\texttt{2.  The counterfactual condition: The individual is not Jyka (\$Jyka = 0\$).} \\
\\ 
\texttt{We are given the following probabilities based on the world's conditions:} \\
\texttt{*   \$P(Lirg | Jyka = 1) = 49\textbackslash{}\%\$} \\
\texttt{*   \$P(Lirg | Jyka = 0) = 63\textbackslash{}\%\$} \\
\\ 
\texttt{The question effectively asks: Is \$P(Lirg | Jyka = 0) > P(Lirg | Jyka = 1)\$?} \\
\\ 
\texttt{Substituting the values:} \\
\texttt{Is \$63\textbackslash{}\% > 49\textbackslash{}\%\$?} \\
\\ 
\texttt{The answer is **Yes**, the probability of seeing \$Lirg\$ is higher when the individual is not Jyka.} \\
\\ 
\texttt{However, the prompt asks to start the answer with "Yes" or "No" followed by reasoning. There is a critical logical trap in the phrasing of the question itself versus the standard interpretation of causal effects. Let's re{-}read the specific question carefully: *"For those who are jyka, would it be more likely to see lirg if the individual was not jyka?"*} \\
\\ 
\texttt{*   **Subject:** An individual who is Jyka.} \\
\texttt{*   **Actual Outcome:** They see Lirg with probability 49\%.} \\
\texttt{*   **Hypothetical Outcome:** If they were *not* Jyka, they would see Lirg with probability 63\%.} \\
\texttt{*   **Comparison:** Is the hypothetical probability (63\%) greater than the actual probability (49\%)?} \\
\\ 
\texttt{Mathematically, \$63 > 49\$. Therefore, in the counterfactual scenario where the individual is not Jyka, it is indeed more likely to see Lirg than in the actual scenario where they are Jyka.} \\
\\ 
\texttt{**Conclusion:**} \\
\texttt{**Yes**} \\
\texttt{If an individual is Jyka, the probability of observing Lirg is 49\%. If that same individual were not Jyka, the probability of observing Lirg would be 63\%. Since 63\% is greater than 49\%, it would be more likely to see Lirg in the counterfactual state (not being Jyka) than in the actual state (being Jyka). This demonstrates a negative causal effect of Jyka on Lirg in this specific hypothetical world.} \\
            }
        }
    }
    & & \\ \\

    \theutterance \stepcounter{utterance}  
    & & & \multicolumn{2}{p{0.3\linewidth}}{
        \cellcolor[rgb]{0.9,0.9,0.9}{
            \makecell[{{p{\linewidth}}}]{
                \texttt{\tiny{[GM$|$GM]}}
                \texttt{No} \\
            }
        }
    }
    & & \\ \\

    \theutterance \stepcounter{utterance}  
    & & & \multicolumn{2}{p{0.3\linewidth}}{
        \cellcolor[rgb]{0.9,0.9,0.9}{
            \makecell[{{p{\linewidth}}}]{
                \texttt{\tiny{[GM$|$GM]}}
                \texttt{yes} \\
            }
        }
    }
    & & \\ \\

    \theutterance \stepcounter{utterance}  
    & & & \multicolumn{2}{p{0.3\linewidth}}{
        \cellcolor[rgb]{0.9,0.9,0.9}{
            \makecell[{{p{\linewidth}}}]{
                \texttt{\tiny{[GM$|$GM]}}
                \texttt{game\_result = LOSE} \\
            }
        }
    }
    & & \\ \\

\end{supertabular}
}

\end{document}
